%% 3. 結果（実験 1・2）
\setband{\T{\S3 結果}{\S3 Results}}

% セルの値：#1 = famone|famthree，#2 = コーパス，#3 = 生徒，#4 = K（zero|five|twofive）
\newcommand{\fc}[4]{\N{#1#2#3#4}\,{\color{mMuted}[\N{#1lo#2#3#4}, \N{#1hi#2#3#4}]}}

% 族の図：#1 = fam1|fam3，#2 = ymin，#3 = ymax，#4 = ylabel
\newcommand{\fampanel}[3]{%
  \nextgroupplot[title=#2]
  \addplot[gray!60, no marks, forget plot] coordinates {(0.5,0) (4.5,0)};
  \addplot[sK0] table[x expr=\thisrow{xi}-0.25, y=effect, y error plus expr=\thisrow{hi}-\thisrow{effect}, y error minus expr=\thisrow{effect}-\thisrow{lo}]{#1_#2_K0.csv};
  \addplot[sK5] table[x expr=\thisrow{xi}, y=effect, y error plus expr=\thisrow{hi}-\thisrow{effect}, y error minus expr=\thisrow{effect}-\thisrow{lo}]{#1_#2_K5.csv};
  \addplot[sK25] table[x expr=\thisrow{xi}+0.25, y=effect, y error plus expr=\thisrow{hi}-\thisrow{effect}, y error minus expr=\thisrow{effect}-\thisrow{lo}]{#1_#2_K25.csv};
  #3}
\newcommand{\famchart}[4]{%
\begin{tikzpicture}
\begin{groupplot}[group style={group size=4 by 1, horizontal sep=0.3cm, ylabels at=edge left, yticklabels at=edge left},
  width=0.255\textwidth, height=4.3cm, xmin=0.5, xmax=4.5, ymin=#2, ymax=#3, xtick={1,2,3,4}, xticklabels={A,V,T,AVT},
  ylabel={#4}, ylabel style={font=\tiny}, ymajorgrids, title style={font=\scriptsize, yshift=-4pt}, legend style={at={(0.03,0.97)}, anchor=north west, font=\tiny, cells={anchor=west}}]
\fampanel{#1}{AFEW}{\legend{K = 0, K = 5, K = 25}}
\fampanel{#1}{IEMOCAP}{}
\fampanel{#1}{MELD}{}
\fampanel{#1}{MOSEI}{}
\end{groupplot}
\end{tikzpicture}}

% 実験 2 の図：#1 = K，#2 = xmin，#3 = xmax
\newcommand{\daipanel}[2]{%
  \nextgroupplot[title=#2]
  \addplot[gray!60, no marks, forget plot] coordinates {(0,0.3) (0,10.7)};
  \addplot[xbar, bar width=4.2pt, fill=cProp!70, draw=cProp, error bars/.cd, x dir=both, x explicit, error bar style={mDarkTeal, line width=0.6pt}]
    table[x=effect, y=yi, x error plus expr=\thisrow{hi}-\thisrow{effect}, x error minus expr=\thisrow{effect}-\thisrow{lo}]{dai2_#2_K#1.csv};}
\newcommand{\daichart}[4]{%
\begin{tikzpicture}
\begin{groupplot}[group style={group size=4 by 1, horizontal sep=0.3cm, yticklabels at=edge left},
  width=0.255\textwidth, height=#4, xmin=#2, xmax=#3, ymin=0.3, ymax=10.7, ytick={1,...,10},
  yticklabels={\T{自己学習}{self-training}, \T{エントロピー最小化}{entropy min.}, IM, T3A, \T{ラベル伝播}{label prop.}, \T{自分の予測}{own prediction}, S+T, CORAL, DANN, MME},
  y dir=reverse, xmajorgrids, title style={font=\scriptsize, yshift=-4pt}, tick label style={font=\tiny}]
\daipanel{#1}{AFEW}
\daipanel{#1}{IEMOCAP}
\daipanel{#1}{MELD}
\daipanel{#1}{MOSEI}
\end{groupplot}
\end{tikzpicture}}

% 得の棒グラフ：#1 = fam1，#2 = コーパス
\newcommand{\gainpanel}[2]{%
  \nextgroupplot[title=#2]
  \addplot[gray!60, no marks, forget plot] coordinates {(0.5,0) (4.5,0)};
  \addplot[ybar, bar shift=0pt, bar width=8pt, fill=cK5, draw=cK5, error bars/.cd, y dir=both, y explicit, error bar style={mDarkTeal, line width=0.6pt}]
    table[x expr=\thisrow{xi}-0.19, y=effect, y error plus expr=\thisrow{hi}-\thisrow{effect}, y error minus expr=\thisrow{effect}-\thisrow{lo}]{#1_#2_K5.csv};
  \addplot[ybar, bar shift=0pt, bar width=8pt, fill=cK25, draw=cK25, error bars/.cd, y dir=both, y explicit, error bar style={mDarkTeal, line width=0.6pt}]
    table[x expr=\thisrow{xi}+0.19, y=effect, y error plus expr=\thisrow{hi}-\thisrow{effect}, y error minus expr=\thisrow{effect}-\thisrow{lo}]{#1_#2_K25.csv};}
\newcommand{\gainchart}[1]{%
\begin{tikzpicture}
\begin{groupplot}[group style={group size=4 by 1, horizontal sep=0.35cm, ylabels at=edge left, yticklabels at=edge left},
  width=0.235\textwidth, height=4.2cm, xmin=0.5, xmax=4.5, ymin=-6, ymax=12, xtick={1,2,3,4}, xticklabels={A,V,T,AVT},
  ylabel={\T{得（ポイント）}{gain (points)}}, ylabel style={font=\tiny}, ymajorgrids, title style={font=\scriptsize, yshift=-4pt},
  legend style={at={(0.03,0.97)}, anchor=north west, font=\tiny, cells={anchor=west}}]
\gainpanel{#1}{AFEW}\legend{K = 5, K = 25}
\gainpanel{#1}{IEMOCAP}
\gainpanel{#1}{MELD}
\gainpanel{#1}{MOSEI}
\end{groupplot}
\end{tikzpicture}}

% 絶対の UAR の棒グラフ：#1 = 生徒の記号（A|V|T|AVT），#2 = K，#3 = 題
\newcommand{\levelpanel}[3]{%
  \nextgroupplot[title={#3}]
  \addplot[ybar, bar shift=0pt, bar width=11pt, fill=cNone!60, draw=cNone] table[x expr=\thisrow{ci}-0.2, y=none#2]{levels_#1.csv};
  \addplot[ybar, bar shift=0pt, bar width=11pt, fill=cProp, draw=cProp] table[x expr=\thisrow{ci}+0.2, y=prop#2]{levels_#1.csv};}
\newcommand{\levelchart}[3]{%
\begin{tikzpicture}
\begin{groupplot}[group style={group size=2 by 1, horizontal sep=0.6cm, ylabels at=edge left, yticklabels at=edge left}, width=0.5\textwidth, height=4.1cm,
  xmin=0.4, xmax=4.6, ymin=#2, ymax=#3, xtick={1,2,3,4}, xticklabels={AFEW, IEMOCAP, MELD, MOSEI}, tick label style={font=\scriptsize},
  ylabel={\T{テスト UAR（\%）}{test UAR (\%)}}, ylabel style={font=\scriptsize}, ymajorgrids, title style={font=\small, yshift=-3pt},
  nodes near coords={\pgfmathprintnumber[fixed, precision=1]{\pgfplotspointmeta}}, every node near coord/.append style={font=\tiny},
  ]
\levelpanel{#1}{5}{K = 5\T{（感情ごとに 5 本）}{ (5 labels per emotion)}}
\levelpanel{#1}{25}{K = 25\T{（感情ごとに 25 本）}{ (25 labels per emotion)}}
\end{groupplot}
\end{tikzpicture}}
% 凡例（手で描く）
\newcommand{\levellegend}{{\scriptsize\tikz[baseline=-0.3ex]\fill[cNone!60, draw=cNone] (0,-0.1) rectangle (0.4,0.15);\ \T{K 本だけ}{K labels only}\qquad \tikz[baseline=-0.3ex]\fill[cProp, draw=cProp] (0,-0.1) rectangle (0.4,0.15);\ \T{提案（K 本 ＋ 教師の疑似ラベル）}{proposal (K labels + teacher pseudo-labels)}}}
% 得の表：#1 = 生徒の記号
\newcommand{\gaintable}[1]{\scriptsize\setlength{\tabcolsep}{5pt}\begin{tabular}{@{}lrrrr@{}}\toprule
\T{得（提案 − K 本だけ）[95\% 区間]}{Gain (proposal − K-only) [95\% interval]} & AFEW & IEMOCAP & MELD & MOSEI \\ \midrule
K = 5 & \fc{famone}{AFEW}{#1}{five} & \fc{famone}{IEMOCAP}{#1}{five} & \fc{famone}{MELD}{#1}{five} & \fc{famone}{MOSEI}{#1}{five} \\
K = 25 & \fc{famone}{AFEW}{#1}{twofive} & \fc{famone}{IEMOCAP}{#1}{twofive} & \fc{famone}{MELD}{#1}{twofive} & \fc{famone}{MOSEI}{#1}{twofive} \\ \bottomrule\end{tabular}}

\begin{frame}{\T{結果 A-1（音声の識別器）：別のモダリティの予測を足すと，4 コーパスで上がる}{Result A-1 (audio classifier): the other modalities' predictions lift it in all 4 corpora}}
\rl{\T{実験 1．比べるもの：提案（K 本 ＋ 教師の疑似ラベル）対 K 本だけ．同じ K 本・同じ出発点}{Exp.~1. Compared: proposal (K labels + teacher pseudo-labels) vs K labels only; same labels, same start}}
\hosoku{appfam1}{\T{全セルの数値表（K = 0 を含む）}{full table of cells (incl.\ K = 0)}}
\centering
\levellegend\par\vspace{-0.3em}
\levelchart{A}{30}{85}\par
\vspace{-0.1em}
\footnotesize
\gaintable{A}\par
\takeaway{\T{音声では 12 セルすべてで正，10 セルで区間が 0 より上．MOSEI は検定でも有意（Holm 後 p < 0.001）}{Audio: positive in all 12 cells, interval above 0 in 10; MOSEI also significant (Holm p < 0.001)}}
\end{frame}

\begin{frame}{\T{結果 A-2（言葉の識別器）：IEMOCAP だけ大きく上がる．ほかは揺れの中}{Result A-2 (text classifier): a large gain only in IEMOCAP; elsewhere within the noise}}
\rl{\T{実験 1．比べるもの：提案 対 K 本だけ}{Exp.~1. Compared: proposal vs K labels only}}
\centering
\levellegend\par\vspace{-0.3em}
\levelchart{T}{30}{72}\par
\vspace{-0.1em}
\footnotesize
\gaintable{T}\par
\takeaway{\T{IEMOCAP の言葉は +5.0〜+7.5．AFEW・MELD・MOSEI は −1.2〜+1.6 で区間が 0 をまたぐ}{IEMOCAP text: +5.0 to +7.5. AFEW, MELD, MOSEI: −1.2 to +1.6, intervals crossing 0}}
\end{frame}

\begin{frame}{\T{結果 A-3（顔の識別器）：上がらない．IEMOCAP では下がる}{Result A-3 (face classifier): no gain; a loss in IEMOCAP}}
\rl{\T{実験 1．比べるもの：提案 対 K 本だけ}{Exp.~1. Compared: proposal vs K labels only}}
\centering
\levellegend\par\vspace{-0.3em}
\levelchart{V}{30}{90}\par
\vspace{-0.1em}
\footnotesize
\gaintable{V}\par
\takeaway{\T{顔で区間が 0 より上なのは MOSEI の K = 25 だけ．IEMOCAP は負（理由は \S4：切片が悪くなる）}{Face: the interval is above 0 only for MOSEI at K = 25; IEMOCAP is negative at all K (why: \S4, the intercept gets worse)}}
\end{frame}

\begin{frame}{\T{結果 A-4（AVT：3 つをつないだ識別器）：MELD 以外は上がらない}{Result A-4 (AVT, the concatenated classifier): no gain except in MELD}}
\rl{\T{実験 1．比べるもの：提案 対 K 本だけ}{Exp.~1. Compared: proposal vs K labels only}}
\centering
\levellegend\par\vspace{-0.3em}
\levelchart{AVT}{35}{92}\par
\vspace{-0.1em}
\footnotesize
\gaintable{AVT}\par
\takeaway{\T{AVT は MELD で +1.4〜+3.4，ほかは −0.8〜+1.7 で区間が 0 をまたぐ（理由は \S4：教師が AVT より当たっていない）}{AVT: +1.4 to +3.4 in MELD, −0.8 to +1.7 elsewhere with intervals crossing 0 (why: \S4, the teacher is no better than AVT)}}
\end{frame}

% 実験 2 の表の 1 行：#1 = 手法の記号（daitwo の key），#2 = 手法名
\newcommand{\dairow}[2]{#2 & \N{daitwocellAFEW#1five} & \N{daitwocellIEMOCAP#1five} & \N{daitwocellMELD#1five} & \N{daitwocellMOSEI#1five} \\}

\begin{frame}{\T{結果 B（実験 2）：同じラベル無しの発話を使う標準の手法と比べると，提案は多くに勝つが，全部ではない}{Result B (Exp.~2): against standard methods that use the same unlabeled utterances, the proposal beats most but not all}}
\rl{\T{比べるもの：提案 − 標準の手法（AVT の生徒，K = 5）}{Compared: proposal − standard method (AVT student, K = 5)}}
\hosoku{appdai2}{\T{区間つきの図（K = 0・5・25）}{charts with intervals (K = 0, 5, 25)}}
\begin{columns}[T]
\begin{column}{0.52\textwidth}
\scriptsize\setlength{\tabcolsep}{4pt}
\begin{tabular}{@{}lrrrr@{}}
\toprule
\T{手法（使う情報）}{Method (information used)} & AFEW & IEMOCAP & MELD & MOSEI \\
\midrule
\multicolumn{5}{@{}l}{\textit{\T{K 本 ＋ 同じモダリティのラベル無し}{K labels + same-modality unlabeled data}}} \\
\dairow{self}{\T{自己学習}{self-training}}
\dairow{entmin}{\T{エントロピー最小化}{entropy min.}}
\dairow{im}{IM}
\dairow{tthreea}{T3A}
\dairow{lp}{\T{ラベル伝播}{label propagation}}
\dairow{own}{\T{自分の予測を提案の形で}{own prediction}}
\multicolumn{5}{@{}l}{\textit{\T{＋ ソースのデータも使う}{+ source data}}} \\
\dairow{st}{S+T}
\dairow{coral}{CORAL}
\dairow{dann}{DANN}
\dairow{mme}{MME}
\bottomrule
\end{tabular}
\end{column}
\begin{column}{0.48\textwidth}
\scriptsize
\T{%
\begin{itemize}
\item \textbf{何を比べたか}：結果 A と同じ AVT の識別器を，提案（教師の疑似ラベル）ではなく標準の手法で直した場合との差．正なら提案が高い．\textbf{太字}は区間が 0 より上（提案が確かに高い），\textcolor{mBad}{赤}は区間が 0 より下（手法が確かに高い）
\item 提案は多くの手法に勝つ（12 セル中，エントロピー最小化 \N{daitwoaboveentmin}，T3A \N{daitwoabovetthreea}，MME \N{daitwoabovemme}，CORAL \N{daitwoabovecoral}，DANN \N{daitwoabovedann}）
\item 負けるのは MOSEI の CORAL（\N{daitwoMOSEIcoralfive}）と，AFEW の K = 0 の T3A・エントロピー最小化・IM（補足）
\item 使う情報が違うので，優劣でなく「情報の違い」として読む
\end{itemize}
}{%
\begin{itemize}
\item \textbf{What is compared}: the AVT classifier of Result A, adapted by a standard method instead of the proposal. Positive = proposal higher. \textbf{Bold}: interval above 0; \textcolor{mBad}{red}: interval below 0 (method reliably higher)
\item The proposal beats most methods (of 12 cells: entropy min.\ \N{daitwoaboveentmin}, T3A \N{daitwoabovetthreea}, MME \N{daitwoabovemme}, CORAL \N{daitwoabovecoral}, DANN \N{daitwoabovedann})
\item It loses to CORAL in MOSEI (\N{daitwoMOSEIcoralfive}) and, at K = 0 in AFEW, to T3A, entropy min.\ and IM (appendix)
\item Different information is used, so read as information differences, not a ranking
\end{itemize}
}
\end{column}
\end{columns}
{\tiny\color{mMuted}\T{提案 − 手法（UAR ポイント，50 単位の平均）．ラベル伝播は K ≥ 1 だけ．出どころ：figures/dai2\_*.csv}{Proposal − method (UAR points, mean over 50 units). Label propagation exists only for K ≥ 1. Source: figures/dai2\_*.csv}}
\end{frame}
